%───────────────────────────────────────────────────────────────────────────────
% CATPPUCCIN THEME — Versión completa y robusta (sin errores)
% Para template académico de física, matemática, computación.
%───────────────────────────────────────────────────────────────────────────────

% ─── 1. Selección del flavor (descomentar solo una) ──────────────────────────
\usepackage[Latte,styleAll]{catppuccinpalette}   % Tema claro
%\usepackage[Frappe,styleAll]{catppuccinpalette} % Tema oscuro
%\usepackage[Macchiato,styleAll]{catppuccinpalette}% Tema oscuro
%\usepackage[Mocha,styleAll]{catppuccinpalette}  % Tema oscuro

% ─── 2. Configuración MANUAL de claro/oscuro (según el flavor elegido) ───────
% Si elegiste Latte: \darkthemefalse
% Si elegiste Frappé/Macchiato/Mocha: \darkthemetrue
\newif\ifdarktheme
\darkthemefalse   % ← Cambia a \darkthemetrue si usas un flavor oscuro

% ─── 3. Página y texto base ─────────────────────────────────────────────────
\ifdarktheme
  \pagecolor{CtpBase}
  \AtBeginDocument{\color{CtpText}}
\else
  \pagecolor{white}
  \AtBeginDocument{\color{black}}
\fi

% ─── 4. Todos los colores de la guía (disponibles para usar) ─────────────────
% Fondo y superficie
\colorlet{background}{CtpBase}
\colorlet{secondarybg}{CtpMantle}
\colorlet{crustbg}{CtpCrust}
\colorlet{surface0}{CtpSurface0}
\colorlet{surface1}{CtpSurface1}
\colorlet{surface2}{CtpSurface2}
\colorlet{overlay0}{CtpOverlay0}
\colorlet{overlay1}{CtpOverlay1}
\colorlet{overlay2}{CtpOverlay2}

% Tipografía
\colorlet{maintext}{CtpText}
\colorlet{subtext0}{CtpSubtext0}
\colorlet{subtext1}{CtpSubtext1}
\colorlet{subtletext}{CtpOverlay1}

% Mensajes funcionales
\colorlet{success}{CtpGreen}
\colorlet{warning}{CtpYellow}
\colorlet{error}{CtpRed}
\colorlet{info}{CtpTeal}
\colorlet{difficultycolor}{CtpLavender}

% Enlaces
\colorlet{linkcolor}{CtpBlue}
\colorlet{citecolor}{CtpGreen}
\colorlet{urlcolor}{CtpPeach}
\colorlet{visitedlinkcolor}{CtpLavender}
\colorlet{hoverlinkcolor}{CtpSky}

% Sintaxis / código
\colorlet{keywordcolor}{CtpMauve}
\colorlet{stringcolor}{CtpGreen}
\colorlet{commentcolor}{CtpOverlay2}
\colorlet{paramcolor}{CtpMaroon}
\colorlet{builtincolor}{CtpRed}
\colorlet{typecolor}{CtpYellow}
\colorlet{enumvarcolor}{CtpTeal}
\colorlet{macrocolor}{CtpRosewater}
\colorlet{operatorcolor}{CtpSky}
\colorlet{delimcolor}{CtpOverlay2}

% Diferencias (diff) y depuración
\colorlet{diffheader}{CtpBlue}
\colorlet{diffindex}{CtpOverlay2}
\colorlet{difffile}{CtpPink}
\colorlet{diffhunk}{CtpPeach}
\colorlet{diffchangedbg}{CtpBlue!15}       % con opacidad 15%
\colorlet{diffinsertedbg}{CtpGreen!15}
\colorlet{diffremovedbg}{CtpRed!15}
\colorlet{breakpointcolor}{CtpRed}

% Arcoíris (para bloques destacados, corchetes, etc.)
\colorlet{rainbow1}{CtpRed}
\colorlet{rainbow2}{CtpPeach}
\colorlet{rainbow3}{CtpYellow}
\colorlet{rainbow4}{CtpGreen}
\colorlet{rainbow5}{CtpSapphire}
\colorlet{rainbow6}{CtpLavender}

% Colores para títulos (estructura del documento)
\colorlet{chaptercolor}{CtpPink}
\colorlet{sectioncolor}{CtpSky}
\colorlet{subsectioncolor}{CtpLavender}
\colorlet{subsubsectioncolor}{CtpPeach}

% ─── 5. Colores semánticos para entornos (preservados) ───────────────────────
\colorlet{defcolor}{CtpBlue}
\colorlet{thmcolor}{CtpMauve}
\colorlet{propcolor}{CtpGreen}
\colorlet{ejemcolor}{CtpPeach}
\colorlet{obscolor}{CtpOverlay0}
\colorlet{fisicacolor}{CtpSapphire}
\colorlet{estrellacolor}{CtpYellow}
\colorlet{compcolor}{CtpPink}

% ─── 6. Colores "light" (fondos de cajas) ────────────────────────────────────
\newcommand{\createLightColor}[2]{%
  \ifdarktheme
    \colorlet{#1light}{CtpBase!95!#2}%
  \else
    \colorlet{#1light}{white!97!#2}%
  \fi
}
\createLightColor{defcolor}{defcolor}
\createLightColor{thmcolor}{thmcolor}
\createLightColor{propcolor}{propcolor}
\createLightColor{ejemcolor}{ejemcolor}
\createLightColor{obscolor}{obscolor}
\createLightColor{fisicacolor}{fisicacolor}
\createLightColor{estrellacolor}{estrellacolor}
\createLightColor{compcolor}{compcolor}
% Opcional: también para success, warning, etc.
\createLightColor{success}{success}
\createLightColor{warning}{warning}
\createLightColor{error}{error}
\createLightColor{info}{info}

% ─── 7. Atajos de texto ──────────────────────────────────────────────────────
\newcommand{\textsuccess}[1]{\textcolor{success}{#1}}
\newcommand{\textwarning}[1]{\textcolor{warning}{#1}}
\newcommand{\texterror}[1]{\textcolor{error}{#1}}
\newcommand{\textinfo}[1]{\textcolor{info}{#1}}
\newcommand{\textkeyword}[1]{\textcolor{keywordcolor}{#1}}
\newcommand{\textstring}[1]{\textcolor{stringcolor}{#1}}
\newcommand{\textcomment}[1]{\textcolor{commentcolor}{#1}}

% ─── 9. Colores por defecto para listings / tcolorbox (ejemplo) ──────────────
\colorlet{codebg}{surface0}
\colorlet{codeframe}{overlay1}